\section{Introduction}\label{sec:intro}

A polygon is \emph{simple} if its boundary is a Jordan curve. Suppose a simple polygon has $k$ sides
$S_1,\dots,S_k$ and each side $S_j$ lies on a line $L_j$ taken from a fixed set of $n$ lines. How
large can $k$ be? We write $M(n)$ for the maximum, over all sets of $n$ lines and all such polygons.

This question was asked on MathOverflow by Johnathan Will (\url{https://mathoverflow.net/users/596423})
on July 31, 2026~\cite{MO513798}. Their upper bound comes from the following observation. Consecutive
sides lie on distinct lines, and since the vertex $L_j\cap L_{j+1}$ is visited only once, no unordered
pair of lines occurs twice as a consecutive pair. So the cyclic sequence $L_1,\dots,L_k$ is a closed
trail in the complete graph $K_n$, and $k$ is at most the maximum length of a closed trail in $K_n$, which is
$n\lfloor (n-1)/2\rfloor$ (OEIS \href{https://oeis.org/A050187}{A050187}~\cite{A050187}).

The asker conjectured that $M(n)=n\lfloor (n-1)/2\rfloor$ for every $n$. They gave simple examples for
$n\le 6$, using the pentagram outline for $n=5$ and the Star of David for $n=6$, and self-intersecting
examples based on star polygons for $n=7,9,10$; the latter construction breaks down at $n=11$. In the
comments, the user DSM (\url{https://mathoverflow.net/users/12613}) found a simple $24$-gon on $8$ lines
and remarked that they could reach $18$ sides with $7$ lines but not $21$. The asker then asked whether
one can prove that for $n=7$ only a self-intersecting (``genus~1'') polygon can reach $21$ sides. As of
October 8, 2026 the question had no answers.

\paragraph{Results.} Write $n=2r+1$ for odd $n$; then $n\lfloor (n-1)/2\rfloor=\binom n2$.

\begin{theorem}[Parity obstruction]\label{thm:parity}
Let $n=2r+1$. In every simple arrangement of $n$ lines, and more generally of $n$ pseudolines, the
\emph{alternating $2$-factor} $G$ (Definition~\ref{def:G}) has a number of connected components congruent to
$r(r+1)/2=(n^2-1)/8$ modulo $2$.
\end{theorem}

\begin{theorem}[Reduction]\label{thm:reduction}
If $n$ is odd and a simple polygon with $\binom n2$ sides has its sides on $n$ lines, then the
arrangement is simple and the polygon is exactly the alternating $2$-factor, which must therefore be
connected. If instead such a polygon has $k\neq\binom n2$ sides, then $k\le\binom n2-3$.
\end{theorem}

\begin{corollary}\label{cor:main}
If $n\equiv\pm1\pmod 8$, then $M(n)\le\binom n2-3$. In particular $M(7)\le 18$ and $M(9)\le 33$, so
the conjectured equality fails for every $n\equiv\pm 1\pmod 8$.
\end{corollary}

\begin{theorem}[Constructions; Lean-certified]\label{thm:constructions}
There are explicit arrangements of $7$, $9$, $11$ and $13$ lines with integer coefficients carrying
simple polygons with $18$, $33$, $55$ and $78$ sides respectively (\cref{sec:constructions}).
Consequently
\[
  M(7)=18,\qquad M(9)=33,\qquad M(11)=55=11\cdot 5,\qquad M(13)=78=13\cdot 6 .
\]
\end{theorem}

The values $M(11)$ and $M(13)$ follow from the constructions together with the elementary upper bound.
The values $M(7)$ and $M(9)$ also need Corollary~\ref{cor:main}. The value $M(7)=18$ confirms DSM's
experimental $18$ and shows that it is optimal.

The proof of Theorem~\ref{thm:parity} is elementary. A left-to-right sweep of the wiring diagram shows
that every cycle of $G$ contains an odd number of crossings of a certain type (Lemma~\ref{lem:percycle}),
and an explicit potential function on permutations counts those crossings modulo $2$
(Lemma~\ref{lem:A}). Theorem~\ref{thm:parity} is a statement about \emph{every} arrangement, so for
$n\equiv\pm1\pmod 8$ the obstruction cannot be avoided by choosing the lines cleverly.

\paragraph{Relation to arrangement graphs.} The \emph{arrangement graph} of a simple arrangement has
the intersection points as vertices and the segments between consecutive intersection points on a
line as edges~\cite{BoseEverettWismath2003,FelsnerHurtadoNoyStreinu2006}. A simple polygon with
$\binom n2$ sides on $n$ lines is the same thing as a Hamiltonian cycle of the arrangement graph that
\emph{turns} at every vertex, i.e.\ never uses two consecutive edges of the same line
(Remark~\ref{rem:converse}), and by Theorem~\ref{thm:reduction} the only candidate for such a cycle is
the alternating $2$-factor. Bose, Everett and Wismath~\cite{BoseEverettWismath2003} showed that
\emph{some} arrangements of $n$ lines have non-Hamiltonian arrangement graphs, for $n=6$ and for all odd
$n>6$. Corollary~\ref{cor:main} is a statement of a different kind: for $n\equiv\pm1\pmod 8$ \emph{no}
arrangement of $n$ lines or pseudolines has a turning Hamiltonian cycle, whereas for $n=11$ and $n=13$
some do. The related notion of an \emph{inducing polygon} goes back to~\cite{BoseEverettWismath2003}, who
asked whether every simple arrangement of $n$ lines carries a simple $n$-gon with exactly one side on
each line. Scharf and Scherfenberg~\cite{ScharfScherfenberg2011} relaxed the notion to simple polygons
whose sides lie on the lines of the arrangement and which use every line (several sides may lie on the
same line), i.e.\ exactly the polygons considered here, and showed that every simple arrangement has
such a polygon with $O(n)$ sides; Ackerman, Pinchasi, Scharf and
Scherfenberg~\cite{AckermanPinchasiScharfScherfenberg2013} then answered the original question: an
inducing simple $n$-gon always exists. In this language, $M(n)$ is the largest number of sides of an
inducing polygon of an arrangement of at most $n$ lines, whereas the cited work is concerned with
inducing polygons with \emph{few} sides. We are not aware of earlier work on the maximum number of
sides, on the alternating $2$-factor, or on the parity of its number of components.

\paragraph{Verification status.}
\begin{itemize}
\item The four constructions of Theorem~\ref{thm:constructions} are checked in Lean~4.34.1 (core only, no
  Mathlib) by kernel evaluation with exact integer arithmetic. The relevant theorems are
  \texttt{mo513798\_n7}, \texttt{\_n9}, \texttt{\_n11} and \texttt{\_n13}; the only axiom used is
  \texttt{propext}. See \cref{sec:formal}. They were also rechecked, independently of Lean, in exact
  rational arithmetic (\cref{sec:limitations}).
\item Theorems~\ref{thm:parity} and~\ref{thm:reduction} and Corollary~\ref{cor:main} are proved on paper
  and \emph{not} formally verified. Their conclusions were confirmed by exhaustive computation over all
  simple arrangements of $n\le 9$ pseudolines, i.e.\ over all $112\,018\,190$ commutation classes of
  reduced words of the longest permutation for $n=9$, and by random tests up to $n=23$; this is
  evidence, not a proof (\cref{sec:limitations}).
\end{itemize}

\paragraph{Organisation.} \Cref{sec:reduction} proves Theorem~\ref{thm:reduction}. \Cref{sec:parity} proves
Theorem~\ref{thm:parity}. \Cref{sec:constructions} gives the constructions and \cref{sec:formal} their formal
certification. \Cref{sec:limitations} lists what is not verified, the computational checks, and some open
questions.
